\pdfminorversion=7%
\documentclass[aspectratio=169,mathserif,notheorems]{beamer}%
%
\xdef\bookbaseDir{../../bookbase}%
\xdef\sharedDir{../../shared}%
\RequirePackage{\bookbaseDir/styles/slides}%
\RequirePackage{\sharedDir/styles/styles}%
\toggleToGerman%
%
\subtitle{10.~Fabrik-Datenbank:~Tabelle \emph{customer}}%
%
\begin{document}%
%
\startPresentation%
%
\section{Einleitung}%
%
\begin{frame}%
\frametitle{Einleitung}%
\begin{itemize}%
\item Unsere Datenbank hat nun genau eine Tabelle, nämlich~\sqlil{product}.%
\item<2-> Wir wollen aber nicht nur Produktinformationen speichern, sondern auch Kundeninformationen.%
\item<3-> In unserem Beispiel werden wir also auch mit Kundendaten spielen.%
\item<4-> Diese sind hauptsächlich Textdaten.%
\item<5-> Das gibt uns die Gelegenheit, mit einigen fortgeschrittenen Text-Features eines DBMS zu spielen.%
\item<6-> Gleichzeitig werden wir Möglichkeiten kennenlernen, um die Wertebereiche für Spalten weiter einzuschränken.%
\end{itemize}%
\end{frame}%
%
\section{Erstellen der Tabelle}%
%
\begin{frame}[t]%
\frametitle{Grobstruktur}%
\begin{itemize}%
\item Wir nennen unsere Tabelle für Kunden~\textil{customer}.%
\item<2-> Wir wollen folgende Informationen zu jedem Kunden darin speichern\uncover<3->{%
\begin{itemize}%
\item den Namen\uncover<4->{,}%
\item<4-> die Telefonnummer\uncover<5->{, und}%
\item<5-> die Adresse.%
\end{itemize}%
}%
\item<6-> Insgesamt könnten die Daten also wie folgt aussehen:%
\end{itemize}%
\locateGraphic{6}{width=0.86\paperwidth}{graphics/dbStructureCustomer}{0.07}{0.5}%
\end{frame}%
%
\begin{frame}[t]
\frametitle{Erstellen der Tabelle mit \texttt{CREATE TABLE}}%
\begin{itemize}%
\item Wir wissen ja schon, dass man Tabellen mit \sqlil{CREATE TABLE} erstellt.%
\item<4-> Wir haben dieses Kommando via \psql\ an den \postgresql-Server geschickt und das Erstellen der Tabelle hat funktioniert.%
\item<5-> Wir erkennen einige Komponenten des Kommandos wieder.%
\item<6-> Andere, insbesondere der Teil mit \sqlil{CONSTRAINT}, sind neu.%
\item<7-> Arbeiten wir uns durch das Kommando durch!%
\end{itemize}%
%
\gitLoadAndExecSQL{factory:create_table_customer}{}{factory}{create_table_customer.sql}{factory}{boss}{superboss123}%
\listingSQL{2}{factory:create_table_customer}{0.075}{0.18}{0.9}{0.84}
\listingOutput{3}{factory:create_table_customer}{style=tool_style_sql}{0.05}{0.18}{0.9}{0.84}
\listingSQL{4-}{factory:create_table_customer}{0.2}{0.5}{0.6}{0.499}
\end{frame}
%
\begin{frame}%
\frametitle{Die Spalte \texttt{id}}%
\begin{itemize}%
\item Die erste Spalte nennen wir wieder \sqlil{id}.%
\item<2-> Sie hat den Typ \sqlil{INT GENERATED BY DEFAULT AS IDENTITY PRIMARY KEY}.%
\item<3-> In anderen Worten\uncover<4->{:%
\begin{itemize}%
\item ihre Werte sind ganzzahlig~(\sqlil{INT})\uncover<5->{,}%
\item<5-> werden automatisch von einer automatisch inkrementierten Sequenz generiert~(\sqlil{GENERATED BY DEFAULT AS IDENTITY})\uncover<6->{, und}%
\item<6-> stellen den Primärschlüssel~\inEN{primary key} für unsere Tabelle dar.%
\end{itemize}%
}%
\item<7-> Genaugenommen ist die Spalte wieder ein sogenannter Ersatzschlüssel~\inEN{surrogate key}.%
\end{itemize}%
\end{frame}%
%
\begin{frame}%
\frametitle{Die Spalte \texttt{name}}%
\begin{itemize}%
\item Kunden haben Namen.%
\item<2-> Als nächstes definieren wir also eine Spalte~\sqlil{name}.%
\item<3-> Namen sind Text und Namen sind verschieden lang.\uncover<4->{ %
Nehmen wir also wieder \sqlil{VARCHAR(100)} als Datentyp -- 100~Zeichen sollten reichen für vernünftige Namen.}%
\item<5-> Natürlich braucht jeder Kunde einen Namen.\uncover<6->{ %
Die Spalte wird daher mit \sqlil{NOT NULL} annotiert.%
}%
\item<7-> Müssen Namen einmalig sein?\uncover<8->{ %
Nein.\uncover<9->{ %
Es ist völlig OK wenn zwei Kunden den selben Namen haben.\uncover<10->{ %
Wir müssen die Spalte also \emph{nicht} mit \sqlil{UNIQUE} annotieren.%
}}}%
\end{itemize}%
\end{frame}%
%
\begin{frame}%
\frametitle{Die Spalte \texttt{name}:~\texttt{CONSTRAINT}}%
\begin{itemize}%
\item Namen sind also Zeichenketten mit maximal 100 Zeichen und müssen angegeben werden~(\sqlil{NOT NULL}).
\item<2-> Ist das ausreichend, um sicherzustellen, dass die gespeicherten Kundennamen korrekt sind?%
\item<3-> Wenn wir zu unserer Tabelle~\sqlil{product} zurückgehen, stellen wir fest, dass wir problemlos \sqlil{' bla     blab  '}, \sqlil{' '}, oder sogar die leere Zeichenkette~\sqlil{''} als Produktname eingeben können.\uncover<4->{ Nur weglassen können wir den Namen wegen \sqlil{NOT NULL} nicht\dots}%
\item<5-> Was bedeutet eigentlich \emph{korrekt}?%
\item<6-> Unser DBMS kann nicht feststellen, ob der Kunde Herr Bebbo tatsächlich Bebbo heißt.%
\item<7-> Was wir aber verhinden können, ist \DEzB\ das aus Versehen \inQuotes{ Bebbo} statt \inQuotes{Bebbo} eingegeben wird.%
\item<8-> Und wir können leere Zeichenfolgen verbieten.%
\item<9-> Mit Hilfe von Einschränkungen, die als \sqlil{CONSTRAINT} formuliert werden\cite{PGDG:PD:C}.%
\end{itemize}%
\end{frame}%
%
%
\begin{frame}[t]
\frametitle{\texttt{CONSTRAINT}}%
\begin{itemize}%
\item Beim Erstellen einer Tabelle können wir zusätzliche Einschränkungen als \sqlil{CONSTRAINT} angeben.%
\end{itemize}%
\sqlSyntax{2-}{syntax/create_table_constraint_de.sql}{0.15}{0.26}{0.8}{0.72}%
\end{frame}
%
\begin{frame}%
\frametitle{Die Spalte \texttt{name}:~\texttt{CONSTRAINT}}%
\begin{itemize}%
\item Wir wollen die Werte in der Spalte \sqlil{name} so einschränken, dass sie mit einem \inQuotes{Wort-Zeichen}, \DEzB~\inQuotes{A}, \inQuotes{b}, oder~\inQuotes{张}, beginnen und enden müssen.\uncover<2->{ Dazwischen erlauben wir beliebige Zeichen.}%
\item<3-> Man kann natürlich immer noch \inQuotes{sgjw9345 s熊猫fki345Q} als Name eingeben\only<3>{.}\uncover<4->{ {\dots} aber Leerzeichen oder Zahlen am Anfang und Ende sind unmöglich und leere Zeichenketten auch.}%
\item<5-> Bisher haben wir keine Möglichkeit, so eine Einschränkung zu formulieren.%
\item<6-> Mit \sqlil{LIKE} (wie in unsere SQL-Anfrage nach Schuhen\dots) geht es leider nicht.%
\item<7-> Aber mit einem \alert{Regulären Ausdruck}.%
\end{itemize}%
\end{frame}%
%
\begin{frame}%
\frametitle{Reguläre Ausdrücke}%
\begin{itemize}%
\item Reguläre Ausdrücke -- \glsreset{regex}\glspl{regex} -- sind Vorlagen für Zeichenketten, die mit anderen Zeichenketten verglichen werden können.%
\item<2-> Wenn ein regulärer Ausdruck zu einer Zeichenkette passt, dann spricht man von einem \emph{match}.%
\item<3-> \Pglspl{regex} werden von sehr vielen Werkzeugen und Programmiersprachen~(wie \DEzB~\python\cite{programmingWithPython}) unterstützt -- und auch von \sql\ und, daher, auch von \postgresql\cite{PGDG:PD:PRE}.%
\item<4-> \Pglspl{regex} sind ein komplett anderes Fachgebiet.
\item<5-> Sie können darüber in \citetitle{PGDG:PD:PRE}\cite{PGDG:PD:PRE} nachlesen.%
\item<6-> Hier wollen wir sie nur kurz vorstellen und direkt einsetzen.%
\end{itemize}%
\end{frame}%
%
\begin{frame}[t]
\frametitle{Die Spalte \texttt{name}:~\texttt{CONSTRAINT} und Reguläre Ausdrücke}%
\begin{itemize}%
\only<-11>{%
\item Wir spezifieren, dass die Werte in der Spalte~\sqlil{name} immer zum \pgls{regex} \sqlil{'^\\w.*\\w\$'} passen müssen.%
\only<-9>{\item<2-> Das Zeichen~\textil{^} stellt den Beginn des Textes dar.%
\item<3-> \textil{\\w} steht für ein \inQuotes{Wort-Zeichen}.}%
\item<4-> \textil{^\\w} bedeutet also, dass ein Wort-Zeichen am Anfang der Zeichenkette stehen muss.%
\only<-9>{\item<5-> Der Punkt~\textil{.} steht für ein beliebiges Zeichen.%
\item<6-> Der Stern~\textil{*} bedeutet, dass der Ausdruck davor beliebig oft~(0, 1, 2, \dots)~vorkommen darf.}%
\item<7-> \textil{.*} nach dem \textil{^\\w} bedeutet also, dass nach dem anfänglichen Wort-Zeichen beliebig viele andere Zeichen beliebiger Art im Namen vorkommen dürfen.
\only<-9>{\item<8-> Das Dollarzeichen~\textil{\$} passt zum Ende einer Zeichenkette.}%
\item<9-> Wenn wir also schreiben~\textil{\\w\$}, dann bedeutet das, dass am Ende der Zeichenkette wieder ein Wort-Zeichen kommen muss.%
\item<10-> \sqlil{'^\\w.*\\w\$'} passt also zu \inQuotes{Thomas Weise} oder \inQuotes{李白}, \alert{nicht} aber zu \inQuotes{ Bibbo}, \inQuotes{熊猫3}, oder~\inQuotes{李}.}%
\item<11-> Wir schreiben~\sqlil{CONSTRAINT customer_name_ok  CHECK (name ~ '^\\w.*\\w\$')}, wobei \sqlil{name ~ '^\\w.*\\w\$'} bedeutet, dass der Wert für \sqlil{name} dem gegebenen \pgls{regex} entsprechen muss.
\end{itemize}%
\listingSQL{12}{factory:create_table_customer}{0.075}{0.22}{0.9}{0.84}
\end{frame}
%
\begin{frame}[t]
\frametitle{Die Spalte \texttt{phone}}%
\begin{itemize}%
\item Jeder Kunde \alert<5>{muss} eine Telefonnummer haben.\uncover<2->{ Wir nennen die Spalte dafür~\sqlil{phone}.}%
\item<3-> In China haben Hausanschlüsse 10~Ziffern und Mobiltelefonnummern~11\cite{BD2006BDBK:MPNANSBTTMDFMP}.%
\item<4-> Als Datentyp nutzen wir daher \sqlil{VARCHAR(11)}.%
\item<5-> Da jeder Kunde eine eine Telefonnummer haben \alert<5>{muss}, ist die Spalte auch~\sqlil{NOT NULL}.%
\item<6-> Jede Telefonnummer kann auch nur von einem Kunden verwendet werden, also spezifizieren die \sqlil{UNIQUE}-Einschränkung.%
\item<7-> \sqlil{VARCHAR(11)} lässt beliebige Zeichenketten mit maximal 11~Zeichen zu.%
\item<8-> Wir wollen aber nur Telefonnummern zulassen.%
\item<9-> \Pglspl{regex} to the rescue!%
\end{itemize}%
\listingSQL{-2}{factory:create_table_customer}{0.075}{0.22}{0.9}{0.84}
\end{frame}
%
\begin{frame}[t]
\frametitle{Die Spalte \texttt{phone}: \texttt{CONSTRAINT}}%
\begin{itemize}%
\only<-8>{\item Wir wollen nur Telefonnummern zulassen.}%
\item<2-> Wir schreiben daher \sqlil{CONSTRAINT customer_phone_ok CHECK (phone ~ '^\\d\{10,11\}\$')}.%
\only<-8>{\item<3-> Die Werte in der Spalte \sqlil{phone} müssen immer zum regulären Ausdruck \textil{^\\d\{10,11\}\$} passen.%
\item<4-> \textil{\\d}~steht für ein einzelnes Zahlenzeichen, also 0\dots9.%
\item<5-> \textil{\{10,11\}}~bedeutet, dass der Ausdruck davor mindestens 10~Mal und höchstens 11~Mal hintereinander vorkommen darf.%
\item<6-> \textil{^\\d\{10,11\}\$}~bedeutet also, dass direkt am Anfang der Zeichenkette entweder 10 oder 11 Ziffern kommen und danach endet die Zeichenkette.%
\item<7-> Natürlich können wir nicht verhindern, dass ungültige Telefonnummern eingegeben werden.%
\item<8-> Aber wir erzwingen zumindest, dass die eingegebenen Zeichenketten zum generellen Schema chinesischer Telefonnummern passen.}%
\end{itemize}%
\listingSQL{9}{factory:create_table_customer}{0.075}{0.18}{0.9}{0.84}
\end{frame}
%
\begin{frame}[t]
\frametitle{Die Spalte \texttt{address}}%
\begin{itemize}%
\item Fügen wir nun die Spalte \sqlil{address} für die Kundenadresse hinzu.%
\only<-5>{%
\item<2-> Wir nehmen wieder \sqlil{VARCHAR(255)}.%
\item<3-> Wir markieren die Spalte as \sqlil{NOT NULL}.%
\item<4-> Wir markieren sie \emph{nicht} als \sqlil{UNIQUE}.%
\item<5-> Sie können sich selbst ein \sqlil{CONSTRAINT} als Übung ausdenken.%
}%
\item<7-> Damit sind wir mit dem Erstellen der Tabelle~\sqlil{customer} fertig.%
\end{itemize}%
\listingSQL{6}{factory:create_table_customer}{0.075}{0.22}{0.9}{0.84}
\end{frame}
%
\section{Einfügen von Daten}%
%
\begin{frame}[t]%
\frametitle{Daten Einfügen}%
\begin{itemize}%
\item Nun haben wir die Tabelle \sqlil{customer} erstellt, aber sie ist leer.%
\item<2-> Füllen wir sie mit Daten!%
\uncover<-8>{%
\item<3-> Wir haben vier Kunden: Herr~Bibbo, Herr~Bebbo, Frau~Bebba, und Herr~Bobbo\only<-3>{.}\uncover<4->{:%
\begin{itemize}%
\item Herr~Bibbo lebt im Südcampus unserer Universität Hefei~(合肥大学)\uncover<5->{,}%
\item<5-> Herr~Bebbo hat seine Geschäftsadresse im Rathaus der Stadt Chemnitz in Deutschland\uncover<6->{,}%
\item<6-> Frau~Bebba wohnt direkt am Times Square in New York, USA\uncover<7->{, und}%
\item<7-> Herr~Bobbo arbeitet am Eiffelturm in Paris, Frankreich.%
\end{itemize}}%
\item<8-> Ihre Telefonnummern sind 11~Wiederholungen einer einzigen Ziffer.%
}%
\end{itemize}%
%
\locateGraphic{3-}{width=0.8\paperwidth}{graphics/dbStructureCustomer}{0.1}{0.57}%
\end{frame}%
%
%
\begin{frame}[t]
\frametitle{Füllen der Tabelle \texttt{product}}%
\gitLoadAndExecSQL{factory:insert_into_table_customer}{}{factory}{insert_into_table_customer.sql}{factory}{boss}{superboss123}%
\listingSQL{1}{factory:insert_into_table_customer}{0.05}{0.17}{0.9}{0.91}
\listingOutput{2}{factory:insert_into_table_customer}{style=tool_style_sql}{0.05}{0.17}{0.9}{0.91}
\end{frame}
%
\begin{frame}[t]
\frametitle{\texttt{CHECK CONSTRAINT}}%
\begin{itemize}%
\item Nun wollen wir noch ausprobieren, ob die \sqlil{CHECK CONSTRAINTS} wirklich funktioniert.%
\item<2-> Was denken Sie, kann die SQL-Anweisung hier funktioneren? Warum/nicht?%
\item<4-> Sehr schön. Die Integrität unserer Daten wird also geschützt.%
\end{itemize}%
\gitLoadAndExecSQL{factory:insert_into_table_customer_error}{}{factory}{insert_into_table_customer_error.sql}{factory}{boss}{superboss123}%
\listingSQL{2}{factory:insert_into_table_customer_error}{0.05}{0.27}{0.9}{0.72}
\listingOutput{3}{factory:insert_into_table_customer_error}{style=tool_style_sql}{0.05}{0.27}{0.9}{0.72}
\end{frame}
%
\section{Daten Selektieren}%
%
\begin{frame}[t]
\frametitle{Daten Selektieren: \texttt{LIKE}}%
\begin{itemize}%
\only<-3>{%
\item Nun haben wir Daten in die Tabelle \sqlil{customer} eingefügt.%
\item<2-> Holen wir sie wieder heraus!%
}%
\item<3-> Zuersteinmal wählen wir alle Kunden aus, die in Frankreich wohnen.%
\item<4-> \sqlil{...WHERE address LIKE '\%France\%'} limitiert die Ausgabe von \sqlil{SELECT} auf Zeilen, in denen das Wort \inQuotes{France} -- genauso geschrieben -- irgendwo im Wert von \sqlil{address} vorkommt.%
\end{itemize}%
%
\gitLoadAndExecSQL{factory:select_from_table_customer_1}{}{factory}{select_from_table_customer_1.sql}{factory}{boss}{superboss123}%
\listingSQLandOutput{4}{factory:select_from_table_customer_1}{}{0.05}{0.36}{0.9}{0.665}
\end{frame}
%
\begin{frame}[t]
\frametitle{Daten Selektieren: \texttt{ILIKE}}%
\begin{itemize}%
\only<-6>{%
\item Als nächstes wollen wir wissen, wie viele inländische Kunden, also Kunden mit chinesischer Adresse, wir haben -- im Vergleich zur Zahl der ausländischen Kunden.%
}%
\item<2-> Zuerst erstellen wir ein synthetisches Attribute \sqlil{domestic}~\inDE{inländisch}, dass wahr~(\sqlil{TRUE}) ist, wenn ein Kunde in China wohnt und sonst falsch~(\sqlil{FALSE}).%
\only<-6>{%
\item<3-> Das können wir mit \sqlil{address LIKE '\%China\%' AS domestic} machen.%
\item<4-> Wir wollen aber diesmal lieber \sqlil{ILIKE} nehmen: \sqlil{LIKE} beachtet Groß- und Kleinschreibung, \sqlil{ILIKE} nicht.%
\item<5-> Für \sqlil{LIKE} sind \sqlil{'china'} und \sqlil{'China'} verschieden.%
\item<6-> Für \sqlil{ILIKE} sind \sqlil{'china'}, \sqlil{'China'}, und \sqlil{'CHINA'} das selbe.%
}%
\item<7-> \sqlil{t} steht für \sqlil{TRUE} und \sqlil{f} für \sqlil{FALSE}.%
\end{itemize}%
%
\gitLoadAndExecSQL{factory:select_from_table_customer_2}{}{factory}{select_from_table_customer_2.sql}{factory}{boss}{superboss123}%
\listingSQLandOutput{7-}{factory:select_from_table_customer_2}{}{0.05}{0.33}{0.9}{0.665}
\end{frame}
%
\begin{frame}[t]
\frametitle{Daten Selektieren: \texttt{GROUP BY} und \texttt{COUNT}}%
\begin{itemize}%
\item Wir wollen aber wissen, \alert{wie viele} inländische Kunden, also Kunden mit chinesischer Adresse, wir haben -- im Vergleich zur Zahl der ausländischen Kunden.%
\only<-5>{%
\item<2-> Dafür müssen wir die Kunden nach dem Attribut~\sqlil{domestic} gruppieren und dann zählen.%
\item<3-> Das Gruppieren funktioniert, in dem wir \sqlil{GROUP BY domestic} an die Anfrage anhängen.%
\item<4-> Statt die Kundennamen zu drucken, zeigen wir dann \sqlil{COUNT(*)} an.%
\item<5-> Wie \sqlil{AVG} ist \sqlil{COUNT} eine Aggregatfunktion, die eine Statistik über alle Elemente der Gruppe berechnet {\dots} und diese Statistik ist hier einfach \inQuotes{Anzahl}.%
}%
\item<6-> Wir haben also einen inländischen und drei ausländische Kunden.%
\end{itemize}%
%
\gitLoadAndExecSQL{factory:select_from_table_customer_3}{}{factory}{select_from_table_customer_3.sql}{factory}{boss}{superboss123}%
\listingSQLandOutput{6-}{factory:select_from_table_customer_3}{}{0.05}{0.33}{0.9}{0.665}
\end{frame}
%
\begin{frame}%
\frametitle{Adressen sind kompliziert}%
\begin{itemize}%
\item Wenn Sie gut aufgepasst haben, dann merken Sie, dass unsere Art der Adressbehandlung problematisch ist.%
\item<2-> Zum einen erkennnen wir zwar \sqlil{'China'} und \sqlil{'china'} beides als China\only<-2>{\dots}\uncover<3->{, würden 中国 aber nicht erkennen.}%
\item<3-> Zum anderen würden wir eine \emph{Donut Factory Vive la France} in Shanghai, China sowohl als Französisch als auch als inländisch erkennen.%
\item<4-> Um das letztere Problem zu lösen, könnten wir die Adresse in mehrere Spalten aufteilen, mit einer Spalte für das Land.%
\item<5-> Selbst dann hätten wir noch Probleme mit verschiedenen Schreibweisen für Länder, \DEzB~China, china, 中国, People's Republic of China, PRC, P.R.C., 中华人民共和国, République populaire de Chine, Chine, usw.%
\item<6-> Als endgültige Lösung würden wir wohl eine weitere Tabelle nur für Länder brauchen, die von unserer Kundentabelle refereziert wird.%
\item<7-> Das wollen wir jetzt nicht tun {\dots} aber in der nächsten Einheit lernen Sie, wie das prinzipiell geht.%
\end{itemize}%
\end{frame}%
%
\section{Zusammenfassung}%
%
\begin{frame}%
\frametitle{Zusammenfassung}%
\begin{itemize}%
\item Jetzt haben wir unsere zweite Tabelle erstellt, die Tabelle~\sqlil{customer}.%
\item<2-> Wir haben einige neue Funktionen zum Arbeiten mit Texten gelernt, \DEzB~\pglspl{regex} und \sqlil{ILIKE}.%
\item<3-> Wir haben gelernt, wie wir mit \sqlil{CONSTRAINT} zusätzliche, komplexere Einschränkungen beim Erstellen einer Tabelle definieren kann.%
\item<4-> Wir haben gelernt, dass wir Daten in \sqlil{SELECT}-Anfragen mit \sqlil{GROUP BY} gruppieren können und Aggregatfunktionen wir \sqlil{COUNT} über die Gruppen berechnen können.%
\item<5-> Und wir haben ein wenig ein Gefühl dafür entwickelt, dass Daten wie \inQuotes{Adressen} auf den ersten Blick einfach aussehen, auf den zweiten Blick aber doch komplexer seien können als erwartet.%
\end{itemize}%
\end{frame}%
%
\endPresentation%
\end{document}%%
\endinput%
%
